% ---------------------------------------------------------
% Definition: Ring
% ---------------------------------------------------------

\begin{tcolorbox}[colback=defbox, colframe=defborder, arc=2pt,
  left=6pt, right=6pt, top=4pt, bottom=4pt,
  title={\small\textbf{Definition --- Ring}},
  fonttitle=\small\bfseries]
\begin{definition}[Ring]
\label{def:ring}
A \textbf{ring} is a non-empty set $R$ together with two binary
operations $+$ and $\cdot$ on $R$, a unary operation $-$ on $R$,
and two distinguished elements $0, 1 \in R$, satisfying the
following axioms:
\begin{enumerate}
  \item[(1)] $p + (q + r) = (p + q) + r$
             \hfill\textit{(additive associativity)}
  \item[(2)] $p \cdot (q \cdot r) = (p \cdot q) \cdot r$
             \hfill\textit{(multiplicative associativity)}
  \item[(3)] $p + q = q + p$
             \hfill\textit{(additive commutativity)}
  \item[(5)] $p + 0 = p$
             \hfill\textit{(additive identity)}
  \item[(6)] $p \cdot 1 = p$
             \hfill\textit{(multiplicative identity --- unit)}
  \item[(7)] $p + (-p) = 0$
             \hfill\textit{(additive inverse)}
  \item[(8)] $p \cdot (q + r) = p \cdot q + p \cdot r$
             \hfill\textit{(left distributivity)}
  \item[(9)] $(q + r) \cdot p = q \cdot p + r \cdot p$
             \hfill\textit{(right distributivity)}
\end{enumerate}
The numbering follows Halmos and Givant~\cite{HalmosGivant2009}.
A ring satisfying axiom~(6) is called a \textbf{ring with unit}.
A ring satisfying $p \cdot q = q \cdot p$ for all $p, q \in R$ is
called a \textbf{commutative ring}.
\end{definition}
\end{tcolorbox}

\begin{remark*}[Standard quantified statement]
\[
\begin{aligned}
&(R, +, \cdot, -, 0, 1) \text{ is a ring} \\
&\quad\Longleftrightarrow\quad
  \bigl(\forall p,q,r \in R\bigr) \\
&\qquad
  p + (q+r) = (p+q)+r \\
&\qquad\land\quad p \cdot (q \cdot r) = (p \cdot q) \cdot r \\
&\qquad\land\quad p + q = q + p \\
&\qquad\land\quad p + 0 = p \\
&\qquad\land\quad p \cdot 1 = p \\
&\qquad\land\quad p + (-p) = 0 \\
&\qquad\land\quad p \cdot (q + r) = p \cdot q + p \cdot r \\
&\qquad\land\quad (q + r) \cdot p = q \cdot p + r \cdot p.
\end{aligned}
\]
\end{remark*}

\begin{remark*}[Negated quantified statement]
\[
\begin{aligned}
&(R, +, \cdot, -, 0, 1) \text{ is \textbf{not} a ring} \\
&\quad\Longleftrightarrow\quad
  \bigl(\exists p,q,r \in R\bigr) \\
&\qquad
  p + (q+r) \neq (p+q)+r \\
&\qquad\lor\quad p \cdot (q \cdot r) \neq (p \cdot q) \cdot r \\
&\qquad\lor\quad p + q \neq q + p \\
&\qquad\lor\quad p + 0 \neq p \\
&\qquad\lor\quad p \cdot 1 \neq p \\
&\qquad\lor\quad p + (-p) \neq 0 \\
&\qquad\lor\quad p \cdot (q + r) \neq p \cdot q + p \cdot r \\
&\qquad\lor\quad (q + r) \cdot p \neq q \cdot p + r \cdot p.
\end{aligned}
\]
\end{remark*}

\begin{remark*}[Failure mode decomposition]
A structure fails to be a ring at the first axiom it violates.
The eight axioms form two groups:
\[
\underbrace{(1),(3),(5),(7)}_{\text{additive abelian group}}
\qquad
\underbrace{(2),(6)}_{\text{multiplicative semigroup with unit}}
\qquad
\underbrace{(8),(9)}_{\text{distributivity}}
\]
The additive group axioms are the most commonly violated in naive
constructions.  A structure satisfying only (1), (3), (5), (7) is
an abelian group under $+$ but carries no multiplication.
Distributivity is the structural link between the two operations;
its failure destroys the ring property even when both operations
are individually well-behaved.
\end{remark*}

\begin{remark*}[Interpretation]
A ring is an algebraic structure abstracting the arithmetic of the
integers.  The additive structure is always a commutative group.
The multiplicative structure is a monoid (associative with unit)
but need not be commutative.  Distributivity is the axiom that
makes the two operations interact coherently --- it is the
algebraic encoding of the familiar rule that multiplication
distributes over addition.  A ring need not have a unit in the
most general definition; here, following Halmos, a unit is
required, making the multiplicative structure a monoid rather
than a mere semigroup.
\end{remark*}





















